% ---
% id: meat-core/req:reported-key-length
% kind: req
% title: "Reported key length"
% summary: >-
%   The engine reports the maximal key length built from the lower bound h on
%   h_low it uses, the right-hand side of the key-length bound at max{0, h}
%   floored and clamped at 0, as its conservative integer output.
% prominence: standard
% status: accepted
% uses: [def:max-key-length, thm:keylength, cor:keylength-bound]
% mentions: [prop:monotone]
% assumes: []
% refs: []
% provenance: newly-derived
% implements: certiqs_key.meat_core.finitekey.keylength:key_length
% lean:
% ---

The engine never has the exact $\hlow$, only a lower bound $h$ on it that a
solver establishes, and Corollary~\ref{cor:keylength-bound} is where it
obtains its key length from one. The key length is a non-negative integer, both steps that
make it one, the floor and the clamp at $0$, go in the safe direction
(Definition~\ref{def:max-key-length}), and, when $h$ depends on none of
$(\lEC,\eev,\epa)$, Corollary~\ref{cor:keylength-bound}\,(a) carries the
monotonicity of Proposition~\ref{prop:monotone} over to it.

\begin{requirement}\label{req:reported-key-length}
The engine reports as \iface{key_length} the \tref{maximal key length}
built from the lower bound $h\le\hlow$ it uses, that is $\ell_{\max}$ of
Definition~\ref{def:max-key-length} with $\max\{0,h\}$ in place of $\hlow$
(Corollary~\ref{cor:keylength-bound}\,(a)), with its floor and its clamp at
$0$: that is its conservative integer-output rule.
\end{requirement}
